% ============================================================
% THEO-CHIR-STATUS-2: The Chiral-Vacuum Breaking Chain (H_4 -> H_4^+) and the
%   V1-Pinning of the FI-C-9 Emergence Upgrade
%
% OPEN-CHIR-1d-β sub-targets 1d-β-i (breaking-chain group theory) + partial
%   1d-β-iii (no axiom-level pseudoscalar fixes the enantiomorph sign).
%
% A Layer-2.5 provisional STRUCTURAL theorem extending THEO-CHIR-STATUS-1. It does
%   NOT derive FI-C-9 (the current-rigor verdict stays V3); it characterizes the
%   emergence UPGRADE: (i) if a chiral vacuum forms it is the SSB H_4 -> H_4^+
%   (index 2) with a Z_2 pseudoscalar order parameter (the enantiomorph = FI-C-9);
%   (iii, partial) no axiom-level pseudoscalar can fix the sign (the unique
%   primitive pseudoscalar is FI-C-9 itself, MERGE-2), so the value-axis is not
%   'derived' at axiom level => V2 is excluded at axiom level => the capacity=Yes
%   outcome is pinned to V1 (emergent mechanism, contingent sign), reopenable to
%   V2 only by a cross-sector pseudoscalar (1d-β-v).
%
% Patch 0654 (Session 149). Verified by
%   code/verify_status_2_breaking_chain.py (CHECK 1 breaking chain + Z_2 grading;
%   CHECK 2 pseudoscalar/V2-exclusion). Both pass.
%
% CHANGELOG
%   v1.0 (Patch 0654, Session 149, 29 May 2026) — initial registration at Layer
%     2.5 (provisional): the H_4 -> H_4^+ breaking chain + Z_2 pseudoscalar order
%     parameter (1d-β-i); the axiom-level V2-exclusion pinning the upgrade to V1
%     (partial 1d-β-iii). Current-rigor verdict UNCHANGED (V3); this sharpens the
%     UPGRADE target, not the placement.
%   v1.1 (Patch 0656, Session 149, 29 May 2026) — multi-AI review cycle CLOSED 3/3,
%     no falsifier (ChatGPT + Grok + Copilot all CONFIRMED at Layer-2.5; STATUS-2
%     judged the source of the pair's informativeness). CALIBRATION (no verdict
%     change): (A) V2-exclusion reworded 'excluded by the presently registered
%     axiom inventory' (ChatGPT Q7); (B) new Remark flagging the Z_2-order-parameter
%     = FI-C-9 step as a framework interpretation layered on the pure group theory
%     (ChatGPT Q4).
% ============================================================

\documentclass[11pt,letterpaper]{article}
\usepackage[utf8]{inputenc}
\usepackage[margin=1in]{geometry}
\usepackage[T1]{fontenc}
\usepackage{amsmath,amssymb,amsfonts,amsthm}
\usepackage{xcolor}
\usepackage{hyperref}
\usepackage{enumitem}
\usepackage{parskip}

\hypersetup{colorlinks=true,linkcolor=blue,citecolor=blue,urlcolor=blue}

\newtheorem{theorem}{Theorem}[section]
\newtheorem{proposition}[theorem]{Proposition}
\newtheorem{lemma}[theorem]{Lemma}
\newtheorem{corollary}[theorem]{Corollary}
\newtheorem{definition}[theorem]{Definition}
\newtheorem{remark}[theorem]{Remark}

\newcommand{\nhat}{\hat{n}}
\newcommand{\sgn}{\operatorname{sign}}
\newcommand{\Po}{\mathsf{P}}
\newcommand{\Te}{\mathsf{T}}
\newcommand{\Hf}{H_4}
\newcommand{\Hfp}{H_4^{+}}
\newcommand{\Ztwo}{\mathbb{Z}_2}
\newcommand{\jnet}{\vec{\jmath}_{\text{net}}}

\title{\textbf{THEO-CHIR-STATUS-2}\\[4pt]
The Chiral-Vacuum Breaking Chain $\Hf \to \Hfp$\\
and the V1-Pinning of the FI-C-9 Emergence Upgrade\\[6pt]
\large A Layer-2.5 Provisional Structural Theorem\\
OPEN-CHIR-1d-$\beta$ sub-targets 1d-$\beta$-i $+$ partial 1d-$\beta$-iii}
\author{Conscious Point Physics --- Substrate Chirality Arc}
\date{Version 1.1 --- Patch 0656, Session 149 --- 29 May 2026 \\[2pt]\small(v1.0 Patch 0654; multi-AI review cycle CLOSED 3/3 at v1.1, no falsifier; all CONFIRMED)}

\begin{document}
\maketitle

\begin{abstract}
\noindent
THEO-CHIR-STATUS-1 placed the current-rigor status of the one currently-identified chirality
primitive, $\text{FI-C-9} = \sgn(\nhat)$, at \textbf{V3} (irreducible primitive --- not yet derived),
with the achievable upgrade being V1 (emergent mechanism, contingent sign). This theorem sharpens the
\emph{upgrade target} (it does not move the placement: the verdict stays V3). \textbf{(1d-$\beta$-i, the
breaking chain.)} A handedness-bearing substrate vacuum is, structurally, the spontaneous breaking of
the achiral 600-cell isometry group $\Hf$ (order 14400, generated by reflections) to its rotation
subgroup $\Hfp$ (order 7200, index 2): the orientation map $\det:\Hf\to\{\pm1\}$ is a homomorphism
with kernel $\Hfp$, the quotient $\Hf/\Hfp \cong \Ztwo$ is the enantiomorph order parameter, and the
two degenerate chiral vacua are its two cosets (related by any reflection; e.g.\
$\operatorname{diag}(-1,1,1,1)$, MERGE-2). The order parameter \emph{is} the pseudoscalar
$\sgn(\nhat) = \text{FI-C-9}$. \textbf{(Partial 1d-$\beta$-iii, the sign.)} Fixing the value of this
$\Ztwo$ order parameter requires a $\Po$-odd (pseudoscalar) quantity: every $\Po$-even invariant
(Gram data, unsigned volumes, the full distance spectrum) is invariant under the reflection relating
the two vacua and so cannot distinguish them. By MERGE-2 the \emph{unique} primitive pseudoscalar is
FI-C-9 itself; hence no axiom-level pseudoscalar can explicitly fix the sign without circularity.
Therefore the value-axis is \emph{not} ``derived'' at the axiom level: \textbf{V2 is excluded at the
axiom level, and the capacity-$=\textsf{Yes}$ outcome is pinned to V1} --- reopenable to V2 only by a
\emph{cross-sector} pseudoscalar (1d-$\beta$-v, e.g.\ a Standard-Model CP-violating phase). \textbf{Net
sharpening of the goal:} chirality is V3 at current rigor, \emph{and when/if FI-C-9 becomes emergent
(via the deep mechanism 1d-$\beta$-ii) it will be exactly V1 --- emergent mechanism, contingent sign
--- not V2}, absent a cross-sector pseudoscalar. Layer 2.5, provisional; no mechanism is derived (V3
stands). Verified by \texttt{code/verify\_status\_2\_breaking\_chain.py}.
\end{abstract}

\tableofcontents

% ============================================================
\section{Setup: what this sharpens, and what it does not}
\label{sec:setup}
% ============================================================

THEO-CHIR-STATUS-1 established the verdict structure $\{$V1, V2, V3$\}$ (capacity axis: is a
chiral-vacuum mechanism derivable? value axis: derived / spontaneous / free) and placed current rigor
at \textbf{V3}: no chiral-vacuum mechanism is registered (1d-$\beta$-ii open), so the capacity axis is
$\textsf{No}$ and FI-C-9 is the one currently-identified irreducible chirality primitive
(\emph{not yet derived}). This theorem delivers two of the remaining structural sub-targets:

\begin{itemize}[leftmargin=1.5em]
\item \textbf{1d-$\beta$-i} --- the breaking chain a chiral vacuum would instantiate (§\ref{sec:chain});
\item \textbf{partial 1d-$\beta$-iii} --- that no axiom-level pseudoscalar can fix the enantiomorph
sign (§\ref{sec:sign}).
\end{itemize}

\noindent
\textbf{It does not move the placement.} No mechanism is derived here; the capacity axis stays
$\textsf{No}$; the current-rigor verdict stays \textbf{V3}. What this theorem sharpens is the
\emph{upgrade target}: it characterizes what the $\textsf{capacity}=\textsf{Yes}$ outcome would be
(V1, not V2). This converts STATUS-1's ``V3, upgradeable to V1'' into ``V3, and the upgrade is
\emph{exactly} V1.''

% ============================================================
\section{1d-$\beta$-i: the breaking chain $\Hf \to \Hfp$}
\label{sec:chain}
% ============================================================

\begin{lemma}[The achiral group and its rotation subgroup]
\label{lem:groups}
The isometry group of the 600-cell is the Coxeter group $\Hf$, of order $14400$, generated by
reflections (hence containing improper, orientation-reversing elements). The orientation map
$\det:\Hf \to \{\pm1\}$, $g \mapsto \det(g)$, is a group homomorphism; its kernel is the
orientation-preserving \emph{rotation subgroup} $\Hfp$, of order $7200$ and index $2$. Consequently
$\Hf/\Hfp \cong \Ztwo$.
\end{lemma}

\begin{proof}
$\Hf$ is the symmetry group of the regular 600-cell (equivalently the 120-cell), a standard fact
($|\Hf| = 14400 = 120^2$); as a finite Coxeter (reflection) group it contains improper isometries.
$\det$ is multiplicative on $O(4)$, so $\det:\Hf\to\{\pm1\}$ is a homomorphism; it is surjective
because $\Hf$ contains a reflection (e.g.\ $R = \operatorname{diag}(-1,1,1,1)$, $\det R = -1$, which
maps the 600-cell to itself --- MERGE-2 achirality lemma, re-verified
\texttt{verify\_status\_2\_breaking\_chain.py} CHECK 1). Its kernel is the proper subgroup $\Hfp$,
index $2$, order $7200$; the quotient is $\Ztwo$. (The det-grading and closure of representative
proper/improper elements are machine-checked, CHECK 1.)
\end{proof}

\begin{theorem}[The chiral-vacuum breaking chain and its order parameter]
\label{thm:chain}
A handedness-bearing (chiral) substrate vacuum is, structurally, a spontaneous breaking
\[
\Hf \;\longrightarrow\; \Hfp \qquad (\text{index } 2,\ \text{quotient } \Ztwo),
\]
in which the achiral symmetry $\Hf$ of the symmetric vacuum is reduced to the rotation subgroup
$\Hfp$ of the broken vacuum. The order parameter labelling the two degenerate broken vacua is the
$\Ztwo$-valued pseudoscalar $\sgn(\nhat) = \text{FI-C-9}$ (the value of the $\det$-coset on the
realized vacuum); the two vacua are exchanged by any reflection in the nontrivial coset $\Hf\setminus\Hfp$.
\end{theorem}

\begin{proof}
A symmetric (achiral) vacuum is invariant under all of $\Hf$, including the reflections; a chiral
vacuum is invariant only under orientation-preserving isometries, i.e.\ its stabilizer lies in $\Hfp$
(Lemma~\ref{lem:groups}). The symmetry reduction is therefore $\Hf \to \Hfp$. Because the index is
$2$ with quotient $\Ztwo$ (Lemma~\ref{lem:groups}), the broken phase is two-fold degenerate, the two
vacua forming the two $\det$-cosets, exchanged by any reflection $r \in \Hf\setminus\Hfp$. The label
distinguishing them is precisely the $\det$-valued $\Ztwo$ pseudoscalar; in the substrate this is the
enantiomorph $\sgn(\nhat) = \text{FI-C-9}$ (the orientation/handedness of the realized configuration;
the geometry itself is achiral by MERGE-2, so the handedness lives entirely in this order parameter).
\end{proof}

\begin{remark}[The order-parameter identification is a framework step]
\label{rem:interp}
A review calibration (Patch 0656): the group-theoretic content of Theorem~\ref{thm:chain} is that a
chiral vacuum carries a $\Ztwo$ order parameter (the $\det$-coset). The further identification of
that $\Ztwo$ with the substrate enantiomorph $\sgn(\nhat) = \text{FI-C-9}$ is a \emph{framework}
step (it reads the abstract orientation label as the registered substrate pseudoscalar), not a pure
group-theoretic deduction. It is sound within CPP but is logically the one interpretive link in an
otherwise standard-group-theory result; flagged here for honesty.
\end{remark}

\begin{remark}[This is the structure, not the dynamics]
\label{rem:structure}
Theorem~\ref{thm:chain} states \emph{what} would break and the order parameter --- pure group theory.
It does \emph{not} show that the substrate dynamics \emph{do} drive the vacuum into the broken phase;
that is the deep, cross-sector mechanism 1d-$\beta$-ii (F.1 §14.17; OPEN-SM-4 $\leftrightarrow$
SS-corpus), untouched here. Hence the capacity axis stays $\textsf{No}$ and the verdict stays V3
(§\ref{sec:setup}).
\end{remark}

% ============================================================
\section{Partial 1d-$\beta$-iii: no axiom-level pseudoscalar fixes the sign}
\label{sec:sign}
% ============================================================

\begin{lemma}[Only a pseudoscalar distinguishes the two enantiomorphs]
\label{lem:pseudoscalar}
Let the two degenerate vacua be related by a reflection $r$ (Theorem~\ref{thm:chain}). A quantity $Q$
can take different values on the two vacua only if $Q$ is $\Po$-odd (a pseudoscalar). Every
$\Po$-even invariant takes the \emph{same} value on both.
\end{lemma}

\begin{proof}
The two vacua are $r$-images of each other. A $\Po$-even quantity is by definition invariant under
the improper $r$, so it agrees on the two vacua and cannot distinguish them. Only a $\Po$-odd quantity
flips under $r$ and so can carry the distinguishing sign. Concretely (CHECK 2): the signed $4$-volume
$D=\det[v_1;v_2;v_3;v_4]$ of an ordered vertex frame is a pseudoscalar and flips, $D \mapsto \det(r)\,D
= -D$; whereas the Gram matrix $G_{ij}=v_i\!\cdot\!v_j$, the unsigned $|D|$, and the entire pairwise
distance spectrum are $r$-invariant (as $r$ is orthogonal).
\end{proof}

\begin{theorem}[Axiom-level V2-exclusion]
\label{thm:v2excluded}
No quantity available at the level of the registered axiom set $A1$--$A11$ can \emph{explicitly}
fix the value of the enantiomorph order parameter $\sgn(\nhat)$. Hence the value axis of
STATUS-1 is not $\textsf{derived}$ at the axiom level, and verdict V2 is excluded at the axiom level
--- more precisely, \emph{excluded by the presently registered axiom inventory} (the exclusion
inherits MERGE-2's primitive-pseudoscalar inventory, hence MERGE-$\alpha$; review-calibration,
Patch 0656).
\end{theorem}

\begin{proof}
By Lemma~\ref{lem:pseudoscalar}, fixing $\sgn(\nhat)$ requires a $\Po$-odd (pseudoscalar) quantity to
couple to it. By THEO-CHIR-MERGE-2 (uniqueness, conditional on MERGE-$\alpha$; the 600-cell is
achiral so geometry supplies none) the \emph{unique} primitive pseudoscalar in the registered
framework is $\text{FI-C-9} = \sgn(\nhat)$ itself. Therefore the only pseudoscalar available to fix
$\sgn(\nhat)$ is $\sgn(\nhat)$ --- circular, not a derivation. No other axiom-level pseudoscalar
exists; so the sign cannot be explicitly derived at the axiom level. By STATUS-1's verdict structure,
$\textsf{value} \ne \textsf{derived}$ at axiom level $\Rightarrow$ V2 is excluded there.
\end{proof}

\begin{corollary}[The upgrade is pinned to V1]
\label{cor:v1pin}
Combining with THEO-CHIR-STATUS-1: should the capacity axis upgrade to $\textsf{Yes}$ (a chiral-vacuum
mechanism is derived, 1d-$\beta$-ii), the resulting verdict is \textbf{V1} (emergent mechanism,
contingent/free sign), \emph{not} V2 --- at the axiom level. Equivalently, the FI-C-9 emergence
target is exactly the standard spontaneous-symmetry-breaking outcome: a derived mechanism with a
spontaneously selected sign.
\end{corollary}

\begin{proof}
By Theorem~\ref{thm:v2excluded} the value axis is $\textsf{spontaneous}$ or $\textsf{free}$ at axiom
level; by STATUS-1 (Thm. ``V1 upgrade condition'') the $\textsf{capacity}=\textsf{Yes}$ outcome with
value $\in\{\textsf{spontaneous},\textsf{free}\}$ is V1.
\end{proof}

\begin{remark}[The reopening: a cross-sector pseudoscalar]
\label{rem:reopen}
The exclusion is \emph{at the axiom level}, i.e.\ within the present registered framework. If a
pseudoscalar were supplied from another sector --- the natural candidate being a Standard-Model
CP-violating phase (1d-$\beta$-v, the FI-C-9 $\leftrightarrow$ electroweak-parity-violation
cross-link, OPEN-SM-4) --- it could couple to and fix $\sgn(\nhat)$, reopening V2. So the precise
claim is: \textbf{V2 is excluded at the axiom level; the upgrade is V1 unless a cross-sector
pseudoscalar (1d-$\beta$-v) supplies an explicit sign-fixing.} This is a clean falsifier (H3$'$).
\end{remark}

% ============================================================
\section{What is and is not established}
\label{sec:caps}
% ============================================================

\textbf{Established (Layer 2.5, provisional).}
\begin{itemize}[leftmargin=1.5em]
\item \textbf{1d-$\beta$-i:} the chiral-vacuum breaking chain is $\Hf\to\Hfp$ (index 2, quotient
$\Ztwo$); the order parameter is the $\det$-pseudoscalar $\sgn(\nhat)=\text{FI-C-9}$; two degenerate
vacua (Thm~\ref{thm:chain}).
\item \textbf{partial 1d-$\beta$-iii:} no axiom-level pseudoscalar fixes the sign; V2 is excluded at
the axiom level (Thm~\ref{thm:v2excluded}); the upgrade is pinned to \textbf{V1}
(Cor~\ref{cor:v1pin}).
\end{itemize}

\textbf{Not established (honest caps).}
\begin{itemize}[leftmargin=1.5em]
\item \textbf{The placement is unchanged --- still V3.} No chiral-vacuum mechanism is derived
(1d-$\beta$-ii open); the capacity axis stays $\textsf{No}$. This theorem sharpens the \emph{upgrade
target}, not the current status (Rem~\ref{rem:structure}).
\item \textbf{V2-exclusion is axiom-level, not absolute} (Rem~\ref{rem:reopen}): a cross-sector
pseudoscalar (1d-$\beta$-v) would reopen V2.
\item \textbf{Inherits MERGE-2's conditionality:} the uniqueness of FI-C-9 as the primitive
pseudoscalar is conditional on MERGE-$\alpha$ and is a ``currently-identified'' count (G6).
\item \textbf{The deep dynamics are untouched:} 1d-$\beta$-ii (the mechanism that would flip the
capacity axis) and 1d-$\beta$-v (the SM cross-link) remain the deferred cross-sector arc.
\item \textbf{Spatial chirality only:} the $\Te$-arrow $\sgn(\delta)$ status is the parallel
OPEN-CHIR-2a (the $\jnet$-arrow physicalization), not addressed here.
\end{itemize}

% ============================================================
\section{Falsifiers}
\label{sec:falsifiers}
% ============================================================

\begin{itemize}[leftmargin=1.5em]
\item \textbf{(H1$'$)} the 600-cell isometry group is not $\Hf$ / not index-2 over its rotation
subgroup --- breaks Lemma~\ref{lem:groups} and the breaking chain. (Standard group theory;
machine-checked structurally, CHECK 1.)
\item \textbf{(H2$'$)} a $\Po$-even quantity is exhibited that distinguishes the two enantiomorphs ---
breaks Lemma~\ref{lem:pseudoscalar} (would contradict the orthogonality of the relating reflection).
\item \textbf{(H3$'$)} a pseudoscalar independent of FI-C-9 is registered (at axiom level, MERGE-2 G6;
or cross-sector, 1d-$\beta$-v) --- reopens V2 (Rem~\ref{rem:reopen}).
\item \textbf{(H4$'$)} a chiral-vacuum mechanism is derived (1d-$\beta$-ii) --- flips the capacity
axis to $\textsf{Yes}$ and \emph{upgrades V3 $\to$ V1} (the intended payoff; not a refutation).
\end{itemize}

% ============================================================
\section*{References (internal)}
\addcontentsline{toc}{section}{References (internal)}
% ============================================================

\begin{itemize}[leftmargin=1.5em]
\item THEO-CHIR-STATUS-1 (Patch 0653) --- the verdict structure $\{$V1,V2,V3$\}$ and the V3 placement;
this theorem sharpens its upgrade target.
\item THEO-CHIR-MERGE-2 (Patches 0647--0651) --- 600-cell achirality; FI-C-9 the \emph{unique}
primitive pseudoscalar (conditional on MERGE-$\alpha$; ``currently-identified,'' G6).
\item OPEN-CHIR-1b (audit E13) --- the chiral subgroup / icosahedral rotation sense; a component of
the breaking chain.
\item OPEN-CHIR-1d-$\beta$ scope sketch (Patch 0652) --- the i--v decomposition; this is
1d-$\beta$-i $+$ partial 1d-$\beta$-iii.
\item F.1 §14.17 viability ceiling; OPEN-SM-4 $\leftrightarrow$ SS-corpus --- the home of
1d-$\beta$-ii (the mechanism) and 1d-$\beta$-v (the cross-sector pseudoscalar).
\item \texttt{code/verify\_status\_2\_breaking\_chain.py} --- CHECK 1 (breaking chain $+$ $\Ztwo$
grading), CHECK 2 (pseudoscalar / V2-exclusion).
\end{itemize}

\end{document}
